\section{Safety and Privacy Enforcement}
\label{sec:impl-safety-privacy}

The system runs on the developer's machine, modifies code under developer control, and reads from a local knowledge corpus. Three concerns shape the implementation: repairs are LLM-generated and can change behavior; outbound network traffic must never carry source code; and the retrieval corpus must be tamper-evident at load time. This section describes how each concern is enforced. Heavier configuration detail (exact retry policies, HNSW parameters, Ed25519 signing tooling) is deferred to Appendix~\ref{appendix:impl-results}.

\subsection{Repair Safety}
\label{sec:impl-repair-safety}

Repairs surface as VS Code Quick Fixes and are never auto-applied. Each applicable diagnostic exposes two actions: \emph{Apply Secure Fix} (one-click, reversible via Ctrl+Z) and \emph{Preview Secure Fix (diff)}, which opens a side-by-side diff editor backed by an in-memory text-document content provider on the \texttt{code-guardian-diff:} scheme so no temporary files are written and the source file stays unchanged until the developer explicitly applies. Before either action is offered, the pipeline parses the suggested fix with \texttt{@babel/parser} (TypeScript, JSX, and async plugins enabled). Markdown fences and single-backtick wrappers are stripped first; obvious prose openings (``You should\dots'', ``To fix\dots'') are pre-rejected. The result is recorded as the \texttt{autoApplicable} boolean (with a short rejection reason) and aggregated into an auto-applicable rate exposed by the evaluation harness; passing the parse check guarantees the fix could be applied without breaking the build, not that it is semantically correct.

Repair scope is kept minimal: single-line fixes are preferred (e.g.\ \texttt{eval(input)} $\to$ \texttt{JSON.parse(input)}); function-level refactoring is used only when structural change is needed; cross-file changes are not supported and are flagged for manual remediation. The repair prompt carries the detected vulnerability type and CWE, the surrounding code (function body plus imports), and instructions to preserve behavior while emitting only the repaired code. The LLM has no runtime information (variable types, execution paths, test cases), so repairs rest on static pattern recognition. Application is buffer-only --- no commits, no automated deployment. Developers review visually, run local tests if available, and commit through normal version-control and CI/CD channels. Behavioral-equivalence checking, exploit oracles, and full path-coverage test suites are out of scope; the syntactic auto-applicability check is the deterministic, fleet-wide quality bar that complements per-sample manual review.

\subsection{Network Policy and Zero-Egress Gate}
\label{sec:impl-network-policy}

The extension routes every outbound HTTP(S) call through a single chokepoint, \texttt{src/networkPolicy.ts}. The module exports an \texttt{assertEgressAllowed(url)} function that throws an \texttt{EgressBlockedError} unless the URL's hostname is on a fixed loopback allowlist (\texttt{127.0.0.1}, \texttt{::1}, \texttt{localhost}, \texttt{0.0.0.0}) \emph{or} the user has explicitly opted in via the \texttt{codeGuardian.enableExternalDataFetch} setting (default \texttt{false}). The opt-in path exists for the public CWE/CVE/OWASP metadata refresh; it does not relax the policy for any source-code-bearing call. The vulnerability data manager consults the gate before every fetch and additionally treats any cached file as valid when the gate is closed, so a fully offline machine can still load the bundled vulnerability snapshot. Because the gate is a synchronous assertion at the call site, it cannot be bypassed by code that holds a live HTTP client reference.

\subsection{Signed Corpus and Provenance}
\label{sec:impl-corpus-signing}

Retrieval poisoning is the threat that an attacker swaps the on-disk RAG knowledge base for one carrying malicious ``guidance'' that the LLM then trusts as authoritative context. To detect this attack, the RAG corpus is signed with an Ed25519 detached signature over a SHA-256 manifest of every corpus file. The signing tool walks the knowledge directory, computes a SHA-256 for each file, attaches per-file provenance (source URL plus retrieval timestamp), produces a canonical tab-separated payload sorted by file path, signs the payload, and writes \texttt{security-knowledge/corpus-manifest.json}. The private key never ships with the extension; the public key is bundled into the extension package and into the Docker image. At RAG load time the corpus verifier re-derives the manifest hashes from the on-disk content, recomputes the canonical payload, and verifies the Ed25519 signature against the public key. A hash mismatch, a missing file, or a signature failure makes the load abort under the default \texttt{codeGuardian.requireSignedCorpus=true} setting. Because provenance fields are bound into the signed payload, an attacker cannot rewrite a corpus file's ``retrieved from'' string without invalidating the signature.

\subsection{Containerised Harness}
\label{sec:impl-container}

The evaluation harness is shipped as a Docker image that pins \texttt{node:20.19.0-alpine}, installs the locked dependency set with \texttt{npm ci}, copies the harness, the signed corpus, and the public key into the image, and compiles the extension at build time. The private key is excluded by \texttt{.dockerignore} so images cannot accidentally carry signing material. \texttt{docker-compose.yml} brings up Ollama and the harness on a user-defined bridge network with explicit CPU and memory limits (8 CPUs, 16\,GB) and binds the Ollama port to loopback only. This is the containerized-execution clause of the task description's reproducibility triad and the network-isolation arm of the threat model in one stack: the harness has no path to a non-loopback host, and resource measurements are comparable across hosts because the container caps the harness side identically.
